% Part: first-order-logic
% Chapter: sequent-calculus
% Section: proving-things

\documentclass[../../../include/open-logic-section]{subfiles}

\begin{document}

\iftag{FOL}
      {\olfileid[id]{fol}{seq}{pro}}
      {\olfileid[id]{pl}{seq}{pro}}

\olsection{Contoh \usetoken{P}{derivation}}

\begin{ex}
Berikan sebuah $\Log{LK}$-!!{derivation} untuk sekuen $!A \land !B \Sequent !A$.

Kita mulai dengan menuliskan sekuen akhir yang diinginkan pada bagian
paling bawah !!{derivation}.
\begin{prooftree}
\AxiomC{}
\UnaryInf$!A\land !B \fCenter !A$
\end{prooftree}
Selanjutnya, kita perlu menentukan jenis inferensi yang dapat mempunyai
sekuen bawah berbentuk demikian. Inferensi itu mungkin merupakan aturan
struktural, tetapi sebaiknya kita mulai dengan mencari aturan logis.
Satu-satunya penghubung logis yang muncul dalam sekuen bawah adalah
$\land$, jadi kita mencari aturan $\land$; karena simbol $\land$ muncul
dalam anteseden, aturan yang kita cari adalah \LeftR{\land}.
\begin{prooftree}
\AxiomC{}
\RightLabel{\LeftR{\land}}
\UnaryInf$!A\land !B \fCenter !A$
\end{prooftree}
Ada dua kemungkinan sekuen atas bagi inferensi \LeftR{\land}: sekuennya
dapat berupa $!A \Sequent !A$, atau $!B \Sequent !A$. Jelas bahwa $!A
\Sequent !A$ merupakan sekuen awal (dan ini menguntungkan), sedangkan
$!B \Sequent !A$ pada umumnya tidak dapat diturunkan. Kita isikan sekuen
atas tersebut:
\begin{prooftree}
\Axiom$!A \fCenter !A$
\RightLabel{\LeftR{\land}}
\UnaryInf$!A\land !B \fCenter !A$
\end{prooftree}
Sekarang kita mempunyai $\Log{LK}$-!!{derivation} yang benar untuk sekuen
$!A \land !B \Sequent !A$.
\end{ex}

\begin{ex}
Berikan sebuah $\Log{LK}$-!!{derivation} untuk sekuen $\lnot !A \lor !B
\Sequent !A \lif !B$.

Mulailah dengan menuliskan sekuen akhir yang diinginkan pada bagian paling
bawah !!{derivation}.
\begin{prooftree}
\AxiomC{}
\UnaryInf$\lnot !A \lor !B \fCenter !A \lif !B$
\end{prooftree}
Untuk menemukan aturan logis yang dapat menghasilkan sekuen akhir ini,
kita perhatikan penghubung logis di dalamnya: $\lnot$, $\lor$, dan $\lif$.
Untuk sementara, kita hanya memperhatikan $\lor$ dan $\lif$ karena
keduanya merupakan !!{main operator}s dari !!{sentence}s dalam sekuen
akhir, sedangkan $\lnot$ berada di dalam lingkup penghubung lain sehingga
akan kita tangani kemudian. Dengan demikian, pilihan aturan logis untuk
inferensi terakhir adalah aturan \LeftR{\lor} dan aturan \RightR{\lif}.
Sebenarnya kita dapat memilih salah satunya, tetapi mari pilih aturan
\RightR{\lif} (kalaupun tanpa alasan lain, aturan ini memungkinkan kita
menunda pemisahan menjadi dua cabang). Menurut bentuk inferensi
\RightR{\lif} yang dapat menghasilkan sekuen bawah tersebut, bentuknya
harus sebagai berikut:
\begin{prooftree}
\AxiomC{}
\UnaryInf$ !A, \lnot !A \lor !B \fCenter !B $
\RightLabel{\RightR{\lif}} \UnaryInf$ \lnot !A \lor !B \fCenter !A \lif !B $
\end{prooftree}

Jika kita memindahkan $\lnot !A \lor !B$ ke bagian terluar anteseden,
kita dapat menerapkan aturan \LeftR{\lor}. Menurut skemanya, langkah ini
harus terpecah menjadi dua sekuen atas sebagai berikut:
\begin{prooftree}
\AxiomC{}
\UnaryInf$\lnot !A, !A \fCenter !B$
\AxiomC{}
\UnaryInf$!B, !A \fCenter !B$
\RightLabel{\LeftR{\lor}}
\BinaryInf$ \lnot !A \lor !B, !A \fCenter !B $
\RightLabel{\LeftR{\Exchange}}
\UnaryInf$ !A, \lnot !A \lor !B \fCenter !B $
\RightLabel{\RightR{\lif}}
\UnaryInf$ \lnot !A \lor !B \fCenter !A \lif !B $
\end{prooftree}
Ingat bahwa kita sedang berusaha menelusuri ke atas hingga mencapai
sekuen-sekuen awal; tampaknya kita sudah cukup dekat!{} Cabang kanan hanya
berjarak satu pelemahan dan satu pertukaran dari sekuen awal, lalu selesai:
\begin{prooftree}
\AxiomC{}
\UnaryInf$\lnot !A, !A \fCenter !B$
\Axiom$!B \fCenter !B$
\RightLabel{\LeftR{\Weakening}}
\UnaryInf$!A, !B \fCenter !B$
\RightLabel{\LeftR{\Exchange}}
\UnaryInf$!B, !A \fCenter !B$
\RightLabel{\LeftR{\lor}}
\BinaryInf$\lnot !A \lor !B, !A \fCenter !B $
\RightLabel{\LeftR{\Exchange}}
\UnaryInf$ !A, \lnot !A \lor !B \fCenter !B $
\RightLabel{\RightR{\lif}}
\UnaryInf$ \lnot !A \lor !B \fCenter !A \lif !B $
\end{prooftree}

Sekarang, dengan memperhatikan cabang kiri, satu-satunya penghubung logis
dalam sebarang !!{sentence} adalah simbol $\lnot$ pada !!{sentence}s
anteseden, sehingga kita mencari suatu instans aturan \LeftR{\lnot}.
\begin{prooftree}
\AxiomC{}
\UnaryInf$ !A \fCenter !B, !A$
\RightLabel{\LeftR{\lnot}}
\UnaryInf$\lnot !A, !A \fCenter !B$
\Axiom$!B \fCenter !B$
\RightLabel{\LeftR{\Weakening}}
\UnaryInf$!A, !B \fCenter !B$
\RightLabel{\LeftR{\Exchange}}
\UnaryInf$!B, !A \fCenter !B$
\RightLabel{\LeftR{\lor}}
\BinaryInf$\lnot !A \lor !B, !A \fCenter !B $
\RightLabel{\LeftR{\Exchange}}
\UnaryInf$ !A, \lnot !A \lor !B \fCenter !B $
\RightLabel{\RightR{\lif}}
\UnaryInf$ \lnot !A \lor !B \fCenter !A \lif !B $
\end{prooftree}
Seperti cara kita menyelesaikan cabang kanan, kita hanya berjarak satu
pelemahan dan satu pertukaran dari penyelesaian cabang kiri ini.
\begin{prooftree}
\Axiom$!A \fCenter !A$
\RightLabel{\RightR{\Weakening}}
\UnaryInf$ !A \fCenter !A, !B$
\RightLabel{\RightR{\Exchange}}
\UnaryInf$ !A \fCenter !B, !A$
\RightLabel{\LeftR{\lnot}}
\UnaryInf$\lnot !A, !A \fCenter !B$
\Axiom$!B \fCenter !B$
\RightLabel{\LeftR{\Weakening}}
\UnaryInf$!A, !B \fCenter !B$
\RightLabel{\LeftR{\Exchange}}
\UnaryInf$!B, !A \fCenter !B$
\RightLabel{\LeftR{\lor}}
\BinaryInf$\lnot !A \lor !B, !A \fCenter !B $
\RightLabel{\LeftR{\Exchange}}
\UnaryInf$ !A, \lnot !A \lor !B \fCenter !B $
\RightLabel{\RightR{\lif}}
\UnaryInf$ \lnot !A \lor !B \fCenter !A \lif !B $
\end{prooftree}
\end{ex}

\begin{ex}
Berikan sebuah $\Log{LK}$-!!{derivation} untuk sekuen $\lnot !A \lor \lnot !B
\Sequent \lnot (!A \land !B)$.

Dengan teknik di atas, kita mulai dengan menuliskan sekuen akhir yang
diinginkan pada bagian paling bawah.
\begin{prooftree}
\AxiomC{}
\UnaryInf$ \lnot !A \lor \lnot !B \fCenter \lnot (!A \land !B) $
\end{prooftree}
Penghubung utama yang tersedia dari !!{sentence}s dalam sekuen akhir ialah simbol $\lor$
dan simbol $\lnot$. Aturan \LeftR{\lor} maupun \RightR{\lnot} dapat
diterapkan di sini, tetapi kita mulai dengan aturan \RightR{\lnot} karena
untuk sementara aturan itu menghindarkan pemisahan menjadi dua cabang:
\begin{prooftree}
\AxiomC{}
\UnaryInf$!A \land !B, \lnot !A \lor \lnot !B \fCenter $
\RightLabel{\RightR{\lnot}}
\UnaryInf$\lnot !A \lor \lnot !B \fCenter \lnot (!A \land !B)$
\end{prooftree}
Sekarang kita dapat memilih antara aturan \LeftR{\land} dan
\LeftR{\lor}. Mari kita lihat apa yang terjadi jika aturan
\LeftR{\land} diterapkan: kita dapat memulai dengan sekuen $!A,
\lnot !A \lor \lnot !B \Sequent \quad$ atau sekuen $!B, \lnot !A
\lor \lnot !B \Sequent \quad$. Karena !!{derivation} tersebut simetris terhadap
$!A$ dan $!B$, mari kita pilih yang pertama:
\begin{prooftree}
\AxiomC{}
\UnaryInf$!A, \lnot !A \lor \lnot !B \fCenter $
\RightLabel{\LeftR{\land}}
\UnaryInf$!A \land !B, \lnot !A \lor \lnot !B \fCenter $
\RightLabel{\RightR{\lnot}}
\UnaryInf$\lnot !A \lor \lnot !B \fCenter \lnot (!A \land !B)$
\end{prooftree}
Ketika terus mengisi !!{derivation}, kita melihat bahwa kita menghadapi
masalah:
\begin{prooftree}
\Axiom$!A \fCenter !A$
\RightLabel{\LeftR{\lnot}}
\UnaryInf$ \lnot !A, !A \fCenter$
\AxiomC{}
\RightLabel{?}
\UnaryInf$!A \fCenter !B$
\RightLabel{\LeftR{\lnot}}
\UnaryInf$ \lnot !B, !A \fCenter$
\RightLabel{\LeftR{\lor}}
\BinaryInf$\lnot !A \lor \lnot !B, !A \fCenter $
\RightLabel{\LeftR{\Exchange}}
\UnaryInf$!A, \lnot !A \lor \lnot !B \fCenter $
\RightLabel{\LeftR{\land}}
\UnaryInf$!A \land !B, \lnot !A \lor \lnot !B \fCenter $
\RightLabel{\RightR{\lnot}}
\UnaryInf$\lnot !A \lor \lnot !B \fCenter \lnot (!A \land !B)$
\end{prooftree}
Bagian atas cabang kanan tidak dapat direduksi lebih lanjut dan tidak dapat
diubah menjadi sekuen awal melalui inferensi struktural; jadi, ini bukanlah
jalur yang tepat. Dengan demikian, menerapkan aturan \LeftR{\land} di atas
jelas merupakan kekeliruan. Jika kita kembali ke tahap sebelumnya dan
sebagai gantinya menerapkan aturan \LeftR{\lor}, kita memperoleh
\begin{prooftree}
\AxiomC{}
\UnaryInf$\lnot !A, !A \land !B \fCenter $

\AxiomC{}
\UnaryInf$\lnot !B, !A \land !B \fCenter $

\RightLabel{\LeftR{\lor}}
\BinaryInf$\lnot !A \lor \lnot !B, !A \land !B \fCenter $
\RightLabel{\LeftR{\Exchange}}
\UnaryInf$!A \land !B, \lnot !A \lor \lnot !B \fCenter $
\RightLabel{\RightR{\lnot}}
\UnaryInf$\lnot !A \lor \lnot !B \fCenter \lnot (!A \land !B)$
\end{prooftree}
Dengan melengkapi setiap cabang seperti sebelumnya, kita memperoleh
\begin{prooftree}
\Axiom$ !A \fCenter!A$
\RightLabel{\LeftR{\land}}
\UnaryInf$!A \land !B \fCenter !A$
\RightLabel{\LeftR{\lnot}}
\UnaryInf$\lnot !A, !A \land !B \fCenter $

\Axiom$ !B \fCenter !B$
\RightLabel{\LeftR{\land}}
\UnaryInf$!A \land !B \fCenter !B$
\RightLabel{\LeftR{\lnot}}
\UnaryInf$\lnot !B, !A \land !B \fCenter $

\RightLabel{\LeftR{\lor}}
\BinaryInf$\lnot !A \lor \lnot !B, !A \land !B \fCenter $
\RightLabel{\LeftR{\Exchange}}
\UnaryInf$!A \land !B, \lnot !A \lor \lnot !B \fCenter $
\RightLabel{\RightR{\lnot}}
\UnaryInf$\lnot !A \lor \lnot !B \fCenter \lnot (!A \land !B)$
\end{prooftree}
(Dalam langkah-langkah ini kita juga dapat menerapkan aturan $\land$ lebih
bawah daripada aturan $\lnot$ dan tetap memperoleh !!{derivation} yang benar.)
\end{ex}

\begin{ex}
Sejauh ini kita belum menggunakan aturan kontraksi, tetapi aturan itu
kadang-kadang diperlukan. Berikut contoh penggunaannya. Misalkan kita ingin
membuktikan $\quad \Sequent !A \lor \lnot !A$. Menerapkan
$\RightR{\lor}$ secara mundur akan menghasilkan salah satu dari dua
pohon pencarian pembuktian parsial berikut:
\begin{prooftree}
\AxiomC{}
\UnaryInf$ \fCenter !A$
\RightLabel{\RightR{\lor}}
\UnaryInf$ \fCenter !A \lor \lnot !A$
\DisplayProof\qquad\bottomAlignProof
\AxiomC{}
\UnaryInf$!A \fCenter $
\RightLabel{\RightR{\lnot}}
\UnaryInf$ \fCenter \lnot !A$
\RightLabel{\RightR{\lor}}
\UnaryInf$ \fCenter !A \lor \lnot !A$
\end{prooftree}
Tentu saja, tidak satu pun berakhir pada sekuen awal. Kuncinya adalah
menyadari bahwa aturan kontraksi memungkinkan kita menggabungkan dua
salinan !!a{sentence} menjadi satu---dan ketika mencari pembuktian, yakni
bergerak dari bawah ke atas, kita dapat mempertahankan satu salinan $!A
\lor \lnot !A$ dalam premis, misalnya,
\begin{prooftree}
\AxiomC{}
\UnaryInf$ \fCenter !A \lor \lnot !A, !A$
\RightLabel{\RightR{\lor}}
\UnaryInf$ \fCenter !A \lor \lnot !A, !A \lor \lnot !A$
\RightLabel{\RightR{\Contraction}}
\UnaryInf$ \fCenter !A \lor \lnot !A$
\end{prooftree}
Sekarang kita dapat menerapkan $\RightR{\lor}$ untuk kedua kalinya dan
juga memperoleh~$\lnot !A$, yang menghasilkan !!{derivation} lengkap.
\begin{prooftree}
\Axiom$!A \fCenter !A$
\RightLabel{\RightR{\lnot}}
\UnaryInf$\fCenter !A, \lnot !A$
\RightLabel{\RightR{\lor}}
\UnaryInf$\fCenter !A, !A \lor \lnot !A$
\RightLabel{\RightR{\Exchange}}
\UnaryInf$ \fCenter !A \lor \lnot !A, !A$
\RightLabel{\RightR{\lor}}
\UnaryInf$ \fCenter !A \lor \lnot !A, !A \lor \lnot !A$
\RightLabel{\RightR{\Contraction}}
\UnaryInf$ \fCenter !A \lor \lnot !A$
\end{prooftree}
\end{ex}

\begin{prob}
Berikan !!{derivation}s untuk sekuen-sekuen berikut:
\begin{enumerate}
\item $!A \land (!B \land !C) \Sequent (!A \land !B) \land !C$.
\item $!A \lor (!B \lor !C) \Sequent (!A \lor !B) \lor !C$.
\item $!A \lif (!B \lif !C) \Sequent !B \lif (!A \lif !C)$.
\item $!A \Sequent \lnot\lnot !A$.
\end{enumerate}
\end{prob}

\begin{prob}
Berikan !!{derivation}s untuk sekuen-sekuen berikut:
\begin{enumerate}
\item $(!A \lor !B) \lif !C \Sequent !A \lif !C$.
\item $(!A \lif !C) \land (!B \lif !C) \Sequent (!A \lor !B) \lif !C$.
\item $\Sequent \lnot(!A \land \lnot !A)$.
\item $!B \lif !A \Sequent \lnot !A \lif \lnot !B$.
\item $\Sequent (!A \lif \lnot !A) \lif \lnot !A$.
\item $\Sequent \lnot(!A \lif !B) \lif \lnot !B$.
\item $!A \lif !C \Sequent \lnot (!A \land \lnot !C)$.
\item $!A \land \lnot !C \Sequent \lnot (!A \lif !C)$.
\item $!A \lor !B, \lnot !B \Sequent !A$.
\item $\lnot !A \lor \lnot !B \Sequent \lnot(!A \land !B)$.
\item $\Sequent (\lnot !A \land \lnot !B) \lif\lnot(!A \lor !B)$.
\item $\Sequent \lnot(!A \lor !B) \lif (\lnot !A \land \lnot !B)$.
\end{enumerate}
\end{prob}

\begin{prob}
Berikan !!{derivation}s untuk sekuen-sekuen berikut:
\begin{enumerate}
\item $\lnot(!A \lif !B) \Sequent !A$.
\item $\lnot(!A \land !B) \Sequent \lnot !A \lor \lnot !B$.
\item $!A \lif !B \Sequent \lnot !A \lor !B$.
\item $\Sequent \lnot \lnot !A \lif !A$.
\item $!A \lif !B, \lnot !A \lif !B \Sequent !B$.
\item $(!A \land !B) \lif !C \Sequent (!A \lif !C) \lor (!B \lif !C)$.
\item $(!A \lif !B) \lif !A \Sequent !A$.
\item $\Sequent (!A \lif !B) \lor (!B \lif !C)$.
\end{enumerate}
(Semua sekuen ini memerlukan aturan~$\RightR{\Contraction}$.)
\end{prob}
\end{document}
